\section{An outer coloring of the labels}
\label{sec:outer}

The pattern counts allow one coloring of $\cL$ to balance all the
tests needed below. For a fixed cyclic progression as in
\Cref{setup:cyclic}, call a full label \emph{light} if its multiplicity
in the progression is at most $k/M$, and \emph{heavy} otherwise.
A position has the same designation as its label. These designations
always refer to the whole progression, even when we later examine a
short block of its positions.

We first give the probabilistic tool, in its finite asymmetric form.
For the origins of the local lemma, see Erd\H{o}s and
Lov\'asz~\cite{erdosLovasz1975}; the proof of the form used here is included.

\begin{lemma}[Finite asymmetric local lemma]\label{outer:lll}
Let $\mathcal E$ be a finite family of events. Suppose a graph on
$\mathcal E$, with neighbor sets $\Gamma(E)$, has the property that each
$E$ is independent of the sigma-algebra generated by the events outside
$\{E\}\cup\Gamma(E)$. If numbers $z_E\in(0,1)$ satisfy
\begin{equation}\label{outer:lll-condition}
 \Pp(E)\le z_E\prod_{F\in\Gamma(E)}(1-z_F)
 \qquad(E\in\mathcal E),
\end{equation}
then $\Pp(\bigcap_{E\in\mathcal E}E^c)>0$.
\end{lemma}

\begin{proof}
For $T\subseteq\mathcal E$, write $A_T=\bigcap_{F\in T}F^c$.
We prove simultaneously, by induction on $|T|$, that
\begin{equation}\label{outer:conditional}
 \Pp(A_T)>0,
 \qquad
 \Pp(E\mid A_T)\le z_E\quad\text{for every }E\notin T.
\end{equation}
For $T=\varnothing$, positivity is immediate and the conditional bound
follows from \Cref{outer:lll-condition}. Suppose the assertions are
known for smaller sets. Removing any $F\in T$ and applying the
conditional bound gives
\[
 \Pp(A_T)=\Pp(A_{T\setminus\{F\}})
 \bigl(1-\Pp(F\mid A_{T\setminus\{F\}})\bigr)>0.
\]
To prove the conditional bound for $E\notin T$, put
$T_1=T\cap\Gamma(E)$ and $T_2=T\setminus T_1$.
If $T_1=\varnothing$, independence gives
$\Pp(E\mid A_T)=\Pp(E)\le z_E$ directly.
Otherwise, order the members of $T_1$ and condition successively on
their complements, starting with $A_{T_2}$. Each conditioning set has
size at most $|T|-1$, so the induction hypothesis yields
\[
 \Pp(A_{T_1}\mid A_{T_2})
 \ge\prod_{F\in T_1}(1-z_F).
\]
Using independence from $T_2$ in the numerator, we obtain
\[
 \Pp(E\mid A_T)
 \le\frac{\Pp(E\mid A_{T_2})}
             {\Pp(A_{T_1}\mid A_{T_2})}
 \le\frac{\Pp(E)}{\prod_{F\in T_1}(1-z_F)}
 \le z_E.
\]
The last inequality uses \Cref{outer:lll-condition} and
$T_1\subseteq\Gamma(E)$. This completes the induction; take
$T=\mathcal E$ for the conclusion.
\end{proof}

\begin{proposition}[Balanced outer coloring]\label{outer:coloring}
For all sufficiently large $k$, there is a map
$c_*:\cL\to\{0,1\}$ with the following properties.
\begin{enumerate}[label=\textup{(\roman*)}]
 \item In every cyclic progression having at least $\delta k$ light
 positions, each bit occurs on at least one quarter of its light
 positions when labels are colored by $c_*$.
 \item In every eligible recorded pattern of \Cref{count:eligible},
 each bit occurs on at least one quarter of its regular positions.
\end{enumerate}
The threshold on $k$ is absolute.
\end{proposition}

\begin{proof}
Assign to every $\beta\in\cL$ an independent uniform random bit
$\xi_\beta$. We introduce two finite families of bad events.

For the first family, take each distinct full-label word of a cyclic
progression with $L_0\ge\delta k$ light positions. Let $w_{\max}$
be the largest multiplicity of a light label. Order the light positions
by grouping equal labels together, with all choices fixed as functions
of the word. Distribute the resulting list cyclically among
$w_{\max}$ rows. A group has at most $w_{\max}$ positions, so a label
occurs at most once in each row. Every row has at least
\begin{equation}\label{outer:row-size}
 \floor{\frac{L_0}{w_{\max}}}
 \ge\delta M-1
\end{equation}
positions, because $w_{\max}\le k/M$.
For each row let its bad event be that some bit occurs on strictly
fewer than one quarter of its positions. There are at most
$kT_{\mathrm{glob}}$ such events in total, where
$T_{\mathrm{glob}}$ is the number in \Cref{count:global}.

For the second family, take one event for each eligible recorded
pattern and its discrete parameters $h,t$, where $h_0<h\le2M$.
Its bad event is that some bit occurs on strictly fewer than one
quarter of the regular positions. The tested labels are pairwise
distinct by eligibility, and their number is at least $h/2$.
Dummy positions use no variables. Although the realizations of the
patterns have real parameters, these event families are finite:
$\cL$, the word lengths, and the choices of $h,t$ are all finite.

Each event $E$ thus tests $l(E)$ distinct independent label variables.
If $X$ is the number of ones among $l=l(E)$ such variables, then
\[
 \E e^{\pm(X-l/2)}=(\cosh(1/2))^l\le e^{l/8}.
\]
Indeed, $(2j)!\ge2^j j!$ gives
$\cosh t\le e^{t^2/2}$ by termwise comparison of the power series.
Markov's inequality for the two deviations from $l/2$ consequently gives
\begin{equation}\label{outer:tail}
 \Pp(E)
 \le\Pp(X-l/2\le-l/4)+\Pp(X-l/2\ge l/4)
 \le2e^{-l(E)/8}.
\end{equation}
The weak tail inequalities also cover the strict inequalities defining
the bad events, regardless of divisibility of $l$ by four.

Give $E$ the weight $z_E=e^{-l(E)/32}$. By
\Cref{count:global,outer:row-size}, the total weight of all
first-family events is at most
\begin{equation}\label{outer:first-weight}
 kT_{\mathrm{glob}}e^{-(\delta M-1)/32}=o(1).
\end{equation}
For a fixed label $\beta$, \Cref{count:local} bounds the second-family
weight of events using $\xi_\beta$ by
\begin{equation}\label{outer:second-weight}
 \sum_{h_0<h\le2M}(CDh)^{6D}e^{-h/64}
 \le\sum_{h>D^2}(CDh)^{6D}e^{-h/64}=o(1).
\end{equation}
To see the last estimate uniformly, take $C\ge1$. The function
$\log(CDh)/h$ is decreasing for $h\ge D^2$ once $D$ is large,
and hence
\[
 \sup_{h>D^2}\frac{6D\log(CDh)}{h}
 \le\frac{6\log(CD^3)}{D}=o(1).
\]
For all sufficiently large $k$, each summand in the infinite sum in
\Cref{outer:second-weight} is therefore at most $e^{-h/128}$.
Its tail tends to zero as $D\to\infty$.

It follows from \Cref{outer:first-weight,outer:second-weight} that,
uniformly over all labels $\beta$,
\begin{equation}\label{outer:incident-weight}
 \sum_{E:\,\xi_\beta\text{ is used by }E}z_E\le\frac1{500}
\end{equation}
for all sufficiently large $k$. Join two distinct events whenever
their supports share a label variable. Since all label bits are
independent, this is a dependency graph of the kind required in
\Cref{outer:lll}; independence holds from the joint sigma-algebra of
all nonneighbors. By summing \Cref{outer:incident-weight} over the
$l(E)$ labels used by $E$, we have
\[
 \sum_{F\in\Gamma(E)}z_F\le\frac{l(E)}{500}.
\]
The lower bounds $l(E)\ge\delta M-1$ in the first family and
$l(E)\ge h/2>D^2/2$ in the second ensure that $l(E)\ge100$
eventually. In particular $z_E<1/2$. Since
$\log(1-z)\ge-2z$ for $0\le z\le1/2$, we obtain
\begin{align}
 z_E\prod_{F\in\Gamma(E)}(1-z_F)
 &\ge\exp\!\left[-\left(\frac1{32}+\frac2{500}\right)l(E)\right]
 \notag\\
 &\ge2e^{-l(E)/8}\ge\Pp(E).
 \label{outer:criterion}
\end{align}
For the second inequality, the positive difference
$1/8-1/32-2/500$ multiplied by $100$ exceeds $\log2$.
Thus \Cref{outer:lll} gives a choice of all label bits avoiding every
bad event. Fix that choice as $c_*$. Summing the balanced row counts
over the partition of the light positions proves part~\textup{(i)};
the second event family gives part~\textup{(ii)}.
All estimates used absolute constants and depended only on $k$, so
the threshold is absolute.
\end{proof}

Fix such a coloring from now on, and define its pullback to the group by
\begin{equation}\label{outer:pullback}
 c_0(n)=c_*(L(n))\qquad(n\in G).
\end{equation}
